\chapter{Podstawy teoretyczne}
\label{ch:podstawy}

% Jedyny rozdział czysto teoretyczny — do napisania od zera.
% Zbiera pojęcia potrzebne w rozdziałach 3-7.

\section{Weryfikacja dyplomów: stan obecny}
\label{sec:stan-obecny}
% Metody obecne: telefon/e-mail do uczelni, centralne rejestry, PKI.
\TODO{Przegląd obecnych metod weryfikacji i ich ograniczeń.}

\section{Drzewa Merkle}
\label{sec:merkle}
% Definicja, dowód przynależności (proof), złożoność O(log n),
% domenowo separowany liść.
\TODO{Formalne wprowadzenie drzew Merkle i dowodów przynależności.}

\section{Schematy zobowiązań kryptograficznych}
\label{sec:zobowiazania}
% Własności: hiding + binding. Zobowiązanie z solą: c = keccak(wartość, sól).
% Podstawa selektywnego ujawniania (rozdz. 6).
\TODO{Wprowadź zobowiązania hash-owe z solą; własności hiding/binding.}

\section{Rollupy warstwy 2 i EIP-4844}
\label{sec:l2-podstawy}
% Optimistic vs zk rollupy, dostępność danych (DA), blob transactions (blobs),
% podział kosztu: opłata L1-data + wykonanie L2.
\TODO{Model kosztowy rollupów i wpływ EIP-4844 (blobów) na koszt DA.}

\section{Infrastruktura klucza publicznego (PKI) i podpisy cyfrowe}
\label{sec:pki}
% Baza porównawcza dla analizy wartości w rozdz. 7:
% podpisany dokument daje offline'ową, niezależną weryfikację.
\TODO{Opisz PKI jako punkt odniesienia dla analizy wartości blockchaina.}

\section{Ochrona danych osobowych (RODO) i minimalizacja danych}
\label{sec:rodo}
% Zasada minimalizacji danych — motywacja dla selektywnego ujawniania.
\TODO{Zasada minimalizacji danych z RODO jako motywacja prywatności.}
